\chapter{The boundary: the open sector and Conjecture 1}\label{chap:boundary}

\section{The white node}

With Theorem~$7'$ (Chapter~\ref{chap:repair}) the repeated-root sector $p(2,1,1)$ is classified unconditionally; what remains open of the original Theorem~7 is exactly its $p(1,1,1,1)$ clause. Note the nuance: Theorem~$7'$ proves a corrected statement on \texttt{576i2}; the printed bridging claim T9 itself stays refuted (Theorem~\ref{not_T9_bridge_Q}). The gap located in Chapter~\ref{chap:transcript} admits a precise sufficient input: the only
route known to us from the results of \S2.2.4 to the conclusion T0 passes through
Conjecture~1 of Buchholz--MacDougall.

Throughout this chapter, each load-bearing claim carries a label: [K] =
kernel-verified in Lean; [C] = computer-algebra computation (two-system where so
stated; scripts and coverage in \texttt{scripts/CERTIFICATES.md}); [M] = ordinary
mathematical deduction, written but not kernel-checked; [A] = research assessment.
Narrative sentences without a label carry no evidential weight.

\subsection*{The chart, the surface, and the covers}
Fixed notation for the whole chapter.  $B$ denotes the rational $(a,b)$-chart of
the $\sigma_{3}=0$ root space: the quartic has roots $\{1,a,b,c\}$ with
$c=-ab/\mathrm{den}$, $\mathrm{den}=ab+a+b$; its function field is
$K(B)=\Q(a,b)$.  $S$ denotes the \emph{splitting cover} of $B$: the double-double
cover obtained by adjoining $z$ with $z^{2}=R_{q}$ (the $f'$-splitting datum) and
$w$ with $w^{2}=S_{q}$ (the $f''$-splitting datum); its function field is, by
definition, $K(S)=\Q(a,b)(z,w)$ --- the multiquadratic extension.  Every element
of $K(S)$ is a $K(B)$-combination of the spanning set $\{1,z,w,zw\}$ [M: the quadratic tower is
spanned by these four products; it is a basis once Proposition~\ref{prop:S_integral}
is in place].  Kernel statements about $K(S)$ below are proved in span
form --- Theorem~\ref{multiquad_verdicts} quantifies over \emph{any} field
equipped with an embedding of $K(B)$ and square roots of $R_{q}$, $S_{q}$ --- and
reach all of $K(S)$ through that spanning set [M].  Throughout, $\Q(a,b)$ is the base
field of $B$ and $\Q(S)=K(S)$ the splitting extension; the two are never
identified.

The three disjuncts of Theorem~7 and their covers of $S$:

\begin{center}\begin{tabular}{llll}
disjunct & cover & radicand & exclusion \\ \hline
$x_{0}=0$ & $C_{0}\colon v^{2}=\sigma_{1}^{2}-4\sigma_{2}$ & $Q_{4}$ (identity I1) & [K] Thm~\ref{multiquad_verdicts} \\
$x_{0}=\tfrac{3\sigma_{1}+r}{8}$ & $C_{+}\colon v^{2}=s_{1n}(s_{1n}-z)$ & $\gamma_{+}=s_{1n}(s_{1n}-z)$; norm class $-2Q_{4}$ & [K] Thm~\ref{gamma_verdicts} \\
$x_{0}=\tfrac{3\sigma_{1}-r}{8}$ & $C_{-}\colon v^{2}=\gamma_{-}$ & $\gamma_{-}=s_{1n}(s_{1n}+z)$; norm class $-2Q_{4}$ & [K] Thm~\ref{gamma_verdicts} \\
\end{tabular}\end{center}

For $C_{0}$: $\mathrm{den}^{2}(\sigma_{1}^{2}-4\sigma_{2})=Q_{4}$ is the kernel
identity I1, so the $C_{0}$-radicand has square class $Q_{4}$ [K].  For
$C_{\pm}$: the exclusion is a kernel theorem (Theorem~\ref{gamma_verdicts}).
The radicand $\gamma_{\pm}=s_{1n}(s_{1n}\mp\zeta)$ is an element of the layer
$K(\zeta)$, not of the base field, so the biquadratic transfer does not apply to
it directly; instead a layer quadratic step (\texttt{step\_over\_layer}) reduces
a hypothetical square root to a square in the layer or its $S_{q}$-multiple, and
the coordinate norm passage (\texttt{layer\_square\_norm}:
$g_{0}^{2}-g_{1}^{2}R_{q}=(p^{2}-q^{2}R_{q})^{2}$, computed in coordinates, no
Galois action) turns either case, via the kernel identity (I2), into the
statement that $-2Q_{4}$ is a square in $\Q(a,b)$ --- contradicting
\texttt{n2Q4\_not\_square} [K].  The two $C_{\pm}$ radicands are \emph{conjugate}; neither is claimed to have
square class $-2Q_{4}$ itself.  Their product --- equivalently their layer norm ---
has square class $-2Q_{4}$:
$\gamma_{+}\gamma_{-}=s_{1n}^{2}(s_{1n}^{2}-z^{2})=-8s_{1n}^{2}Q_{4}$ by (I2), and
it is this \emph{norm class} that the kernel obstruction consumes.  $C_{0}$ carries
the class $Q_{4}$: three covers, one sextic, two obstruction routes.

\begin{proposition}[Fibre correspondence on the nondegenerate locus]\label{prop:fibre_correspondence}
Let $U\subset S$ be the open locus $\mathrm{den}\neq0$ with the four roots
$\{1,a,b,c\}$ pairwise distinct and nonzero.  At every rational point $P$ of
$U$, each disjunct of Theorem~7 is equivalent to the existence of a rational
point in the corresponding fibre:
$D_{0}(P)\iff\exists\,v\in\Q,\ (P,v)\in C_{0}$, and
$D_{\pm}(P)\iff\exists\,v\in\Q,\ (P,v)\in C_{\pm}$.
\end{proposition}
\begin{proof}[Proof {[M]} with {[K]} identities; every condition listed]
Conditions in force: $\mathrm{den}(P)\neq0$ (chart); roots distinct and nonzero
(nondegeneracy); at a rational point of $S$ the splitting witnesses $z,w$ are
rational numbers, so $\gamma_{\pm}(P)=s_{1n}(s_{1n}\mp z)(P)\in\Q$.
$D_{0}$: the first disjunct asks that $\sigma_{1}^{2}-4\sigma_{2}$ be a rational
square.  The kernel identity (I1),
$\mathrm{den}^{2}(\sigma_{1}^{2}-4\sigma_{2})=Q_{4}$, gives the substitution
$v=\mathrm{den}\cdot t$ between $t^{2}=\sigma_{1}^{2}-4\sigma_{2}$ and
$v^{2}=Q_{4}(P)$; since $\mathrm{den}(P)\neq0$ this is a bijection on solutions
--- clearing the denominator multiplies by the nonzero square
$\mathrm{den}^{2}$ and introduces no spurious solutions.
$D_{\pm}$: the reduction of the disjunct at $x_{0}=(3\sigma_{1}\pm r)/8$ to the
single condition ``$\gamma_{\pm}(P)$ is a rational square'' is the terrain-map
computation (session artifacts, cross-checked [C]; its governing identity (I2)
is kernel); granted the radicand derivation, the fibre statement
$v^{2}=\gamma_{\pm}(P)$ is definitional, the correspondence $v\leftrightarrow-v$
being the only multiplicity.  No further denominators are cleared in this case.
\end{proof}

\begin{lemma}[Squarefree base change]\label{lem:sqfree_basechange}
Let $k$ be a perfect field and $f\in k[a,b]$ nonconstant with a nonconstant
irreducible factor of odd multiplicity.  Then $f$ is not a square in $k'(a,b)$
for any algebraic extension $k'/k$ --- in particular not in $\overline{k}(a,b)$.
\end{lemma}
\begin{proof}[Proof {[M]}]
Over a perfect field irreducible polynomials are separable, hence squarefree;
base change to $k'$ preserves reducedness (geometric reducedness over a perfect
field), so each irreducible factor of $f$ splits over $k'$ into \emph{distinct}
irreducibles, each inheriting the multiplicity it had in $f$.  A nonconstant
factor of odd multiplicity therefore yields a prime of $k'[a,b]$ at which $f$
has odd valuation; a square in $k'(a,b)$ has even valuation at every prime.
\end{proof}

\begin{lemma}[Biquadratic-extension criterion]\label{lem:biquad_criterion}
Let $K$ be a field of characteristic $\neq2$ and $r,s\in K^{\times}$ with $r$,
$s$ and $rs$ all nonsquares in $K$.  Then $K(\sqrt r,\sqrt s)/K$ has degree
four with basis $\{1,\sqrt r,\sqrt s,\sqrt r\sqrt s\}$, and the double-double
cover algebra $K[z,w]/(z^{2}-r,\;w^{2}-s)$ is integral (a field).
\end{lemma}
\begin{proof}[Proof {[M]}]
$[K(\sqrt r):K]=2$ since $r$ is a nonsquare.  If $\sqrt s\in K(\sqrt r)$, write
$\sqrt s=\alpha+\beta\sqrt r$; squaring and comparing coordinates (valid since
$\{1,\sqrt r\}$ is independent and $\operatorname{char}K\neq2$) forces
$2\alpha\beta=0$: $\beta=0$ makes $s=\alpha^{2}$ a square, and $\alpha=0$ makes
$rs=(\beta r)^{2}$ a square --- both excluded.  Hence the tower has degree four
and the product basis; the quotient algebra is then a free $K$-module of rank
four which coincides with the field $K(\sqrt r,\sqrt s)$, hence is integral.
(Standard exponent-2 Kummer theory; see e.g.\ Lang, \emph{Algebra}, Ch.~VI.)
For constant-field extensions the hypotheses persist by
Lemma~\ref{lem:sqfree_basechange}, so the criterion applies verbatim over
$\overline{\Q}(a,b)$.
\end{proof}

\begin{proposition}[Integrality of the splitting cover]\label{prop:S_integral}
$R_{q}$, $S_{q}$ and $R_{q}S_{q}$ are nonsquares in $\Q(a,b)$, and remain
nonsquares in $\overline{\Q}(a,b)$.  Consequently $[K(S):K(B)]=4$, the spanning
set $\{1,z,w,zw\}$ is a basis, and the splitting cover $S$ is geometrically
integral of generic degree four over $B$.
\end{proposition}
\begin{proof}[Proof {[K]}+{[C]}+{[M]}]
Over $\Q$: kernel theorems \texttt{Rq\_not\_square}, \texttt{Sq\_not\_square},
\texttt{RqSq\_not\_square} (\texttt{Gap/BMThm7NormCriterion.lean}), by value
specialisation at witness points ($R_{q}(2,3)=2480$, $S_{q}(1,2)=876$,
$(R_{q}S_{q})(1,2)=171696$, none a rational square) [K].  Over $\overline{\Q}$:
each of the three has a nonconstant factor of odd multiplicity
($R_{q}$ irreducible; $S_{q}=3\cdot(\text{irreducible})$) [C:
\texttt{terrain\_crosscheck}, \texttt{terrain\_class12}], so Lemma~\ref{lem:sqfree_basechange}
applies: none is a square in $\overline{\Q}(a,b)$ ($R_{q}$ and $S_{q}$ each carry
one odd nonconstant factor; they are coprime [C: gcd data], so $R_{q}S_{q}$
retains both).  Degree four, the basis, and integrality now follow from
Lemma~\ref{lem:biquad_criterion} applied over $\Q(a,b)$ and, via
Lemma~\ref{lem:sqfree_basechange}, over $\overline{\Q}(a,b)$ [M].
\end{proof}  In normalised form the conjecture asserts that no
quartic $x(x-1)(x-a)(x-b)$ with four distinct roots is $K$-derived; the unnormalised
conjecture is equivalent to this via the affine group $\langle X^{*}\rangle$, and only
this direction is used here.

\begin{definition}[Conjecture 1, normalised --- OPEN]\label{Conjecture1_normal}\lean{BMThm7Transcript.Conjecture1_normal}
\uses{DistinctRoots, KDerived, bmQuartic}
For every $a,b$ with $\{0,1,a,b\}$ distinct, the quartic $x(x-1)(x-a)(x-b)$ is not
$K$-derived.  This is an open problem; it is stated without proof and carries no
\texttt{leanok} --- it is the white node from which everything below hangs.
\end{definition}

\begin{theorem}[Theorem 7's $p(1,1,1,1)$ clause follows from Conjecture 1]\label{thm7_of_conjecture1}\lean{BMThm7Transcript.thm7_of_conjecture1}\leanok
\uses{Conjecture1_normal, T0_claim}
If Conjecture~1 (normalised) holds, then T0 holds.
\end{theorem}
\begin{proof}\leanok
\uses{Conjecture1_normal}
Vacuously: the conjecture empties the hypothesis class of T0, so the conclusion holds on
an empty domain.  This one line is the only route to T0 known to us (the adjudication record,
P1/P2/P7); what it proves is sufficiency, not necessity.
\end{proof}

\section{Machine-verified algebraic anchors for the terrain map}

The boundary is not merely negative; the terrain surrounding it is mapped in the
$\sigma_{3}=0$ chart of \S2.2.3 (roots $1,a,b,c$ with $c=-ab/\mathrm{den}$,
$\mathrm{den}=ab+a+b$).  The following objects and identities, in the namespace
\texttt{BMThm7Boundary}, are the kernel-certified anchors of the two-system
(Magma~+~Sage) terrain map (RESULTS \S2--10).

\begin{definition}[Chart quantities]\label{boundary-defs}\lean{BMThm7Boundary.Q4}\leanok
$\mathrm{den}=ab+a+b$; $s_{1n}=\sigma_{1}\,\mathrm{den}$; $s_{2n}=\sigma_{2}\,\mathrm{den}$;
$R_{q}=\mathrm{den}^{2}(9\sigma_{1}^{2}-32\sigma_{2})$ (the $f'$-splitting condition,
cleared); $S_{q}=\mathrm{den}^{2}(9\sigma_{1}^{2}-24\sigma_{2})$ (the $f''$-splitting
condition, cleared); the funnel $Q_{4}=(3R_{q}-S_{q})/18=\mathrm{den}^{2}(\sigma_{1}^{2}-4\sigma_{2})$;
and the three pairing curves $F_{1}=ab+a+b-1$, $F_{2}=ab+a-b^{2}+b$, $F_{3}=a^{2}-ab-a-b$.
\end{definition}

\subsection*{The funnel and its factorisation}

All three disjuncts of Theorem~7 are governed by the single sextic $Q_{4}$: the $x_{0}=0$
chart directly, and both $x_{0}=(3\sigma_{1}\pm r)/8$ charts through the norm criterion
via the identity (I2).

\begin{lemma}[I1, cleared]\label{I1_cleared}\lean{BMThm7Boundary.I1_cleared}\leanok
\uses{boundary-defs}
$18\,Q_{4}=3R_{q}-S_{q}$.
\end{lemma}
\begin{proof}\leanok
By \texttt{ring}.
\end{proof}

\begin{lemma}[I2: the C$\pm$ obstruction funnels through $Q_4$]\label{I2}\lean{BMThm7Boundary.I2}\leanok
\uses{boundary-defs}
$s_{1n}^{2}-R_{q}=-8\,Q_{4}$.
\end{lemma}
\begin{proof}\leanok
By \texttt{ring}.
\end{proof}

\begin{lemma}[The factorisation of $Q_4$]\label{Q4_factor}\lean{BMThm7Boundary.Q4_factor}\leanok
\uses{boundary-defs}
$Q_{4}=F_{1}\,F_{2}\,F_{3}$.
\end{lemma}
\begin{proof}\leanok
By \texttt{ring}.
\end{proof}

The kernel statement is the identity only.  From the displayed forms one reads off
that $F_{1}$ is symmetric ($=\mathrm{den}-1$) and $F_{2}(a,b)=-F_{3}(b,a)$; that
$Q_{4}$ is squarefree and that each $F_{i}=0$ is an irreducible plane curve of
genus~$0$ are computer-algebra facts [MC, two-system], not part of the kernel
statement.

\subsection*{The degeneracy locus is BM's own excluded symmetric case}

The pairing identities exhibit, for each root pairing, an exact factorisation of the
symmetry condition with cofactors $-(a+b)$, $(a+1)$, $-(b+1)$.  Away from the vanishing
of these cofactors they identify $Q_{4}=0$ --- the locus where Theorem~7's conclusion is
automatic --- with the locus on which the quartic $\{1,a,b,c\}$ is symmetric under one of
the three root pairings, Buchholz--MacDougall's own excluded case~(v); on the cofactor
loci the identification requires separate elementary inspection.  The kernel content is
the three factorisations themselves.

\begin{lemma}[Pairing $\{1,c\}\mid\{a,b\}$]\label{pairing1}\lean{BMThm7Boundary.pairing1}\leanok
\uses{boundary-defs}
$\mathrm{den}\,(1-a-b)-ab=F_{1}\cdot(-(a+b))$.
\end{lemma}
\begin{proof}\leanok
By \texttt{ring}.
\end{proof}

\begin{lemma}[Pairing $\{1,a\}\mid\{b,c\}$]\label{pairing2}\lean{BMThm7Boundary.pairing2}\leanok
\uses{boundary-defs}
$\mathrm{den}\,(1+a-b)+ab=F_{2}\cdot(a+1)$.
\end{lemma}
\begin{proof}\leanok
By \texttt{ring}.
\end{proof}

\begin{lemma}[Pairing $\{1,b\}\mid\{a,c\}$]\label{pairing3}\lean{BMThm7Boundary.pairing3}\leanok
\uses{boundary-defs}
$\mathrm{den}\,(1-a+b)+ab=F_{3}\cdot(-(b+1))$.
\end{lemma}
\begin{proof}\leanok
By \texttt{ring}.
\end{proof}

\subsection*{Cover nontriviality: the negativity certificates}

A square in the function field $\Q(a,b)$ is non-negative at every rational point where it
is defined.  Each of the eight candidates whose squareness would trivialise a double
cover is \emph{negative} at an explicit rational point, from which the kernel now derives
the verdict that none of the eight is a square --- by the square-class transfer
(Lemma~\ref{sq_transfer}, Theorem~\ref{multiquad_verdicts}) $Q_{4}$ and $-2Q_{4}$
remain non-squares in every field containing $\Q(a,b)$ together with square roots
of $R_{q}$ and $S_{q}$ [K] --- in particular in $K(S)=\Q(a,b)(z,w)$, which \emph{is} the multiquadratic
extension by definition (see the chart/surface/covers section above).
Nontriviality of $C_{0}$ over $K(S)$ is Theorem~\ref{multiquad_verdicts};
nontriviality of $C_{\pm}$ is Theorem~\ref{gamma_verdicts}; geometric
integrality is Proposition~\ref{prop:geomint}.  Point~1 is the $K$-witness $(-51/13,-17/7)$; point~2 is $(1,3)$.  The eight negativity certificates are the kernel anchors; as of 2026-08-02 the
deduction itself is also in the kernel (\texttt{Gap/BMThm7FunctionField.lean}): the
specialisation principle is formalised as evaluation along a ring homomorphism
(Lemma~\ref{spec_principle}), squares descend from the fraction field to the
integrally closed ring $\Q[a][b]$ (Lemma~\ref{sq_descent}), and all eight
non-squareness verdicts in $\Q(a,b)$ --- which is $K(B)$, the function
field of the base chart $B$ (the splitting extension $K(S)$ lies above it) ---
are kernel theorems (Theorem~\ref{covers_nonsquare}).  The two-system [MC] record
remains as the discovery artifact.

\begin{lemma}[Specialisation principle]\label{spec_principle}\lean{BMThm7FunctionField.sq_eval_nonneg}\leanok
A square in $\Q[a][b]$ is non-negative at every rational point.
\end{lemma}
\begin{proof}\leanok
Evaluation is a ring homomorphism; \texttt{map\_pow} and \texttt{sq\_nonneg}.
\end{proof}

\begin{lemma}[Square descent]\label{sq_descent}\lean{BMThm7FunctionField.square_descent}\leanok
$\Q[a][b]$ is a unique factorisation domain, hence integrally closed; an element that
becomes a square in $\Q(a,b)$ is a square in $\Q[a][b]$.
\end{lemma}
\begin{proof}\leanok
A square root in the fraction field is integral (root of $X^{2}-F$), hence lies in the
ring; conclude by injectivity of the localisation map.
\end{proof}

\begin{theorem}[The eight verdicts]\label{covers_nonsquare}\lean{BMThm7FunctionField.Q4_not_square}\leanok
\uses{spec_principle, sq_descent}
None of $Q_{4}$, $Q_{4}R_{q}$, $Q_{4}S_{q}$, $Q_{4}R_{q}S_{q}$, $-2Q_{4}$,
$-2Q_{4}R_{q}$, $-2Q_{4}S_{q}$, $-2Q_{4}R_{q}S_{q}$ is a square in $\Q(a,b)$
(\texttt{Q4\_not\_square} through \texttt{n2Q4RqSq\_not\_square}, eight kernel
theorems consuming the negativity certificates through the evaluation bridges).
\end{theorem}
\begin{proof}\leanok
Square descent, then the specialisation principle at the certificate point
contradicts the kernel negativity.
\end{proof}

\begin{lemma}[Biquadratic square-class transfer]\label{sq_transfer}\lean{BMThm7SquareClass.biquadratic_transfer}\leanok
For fields $K \subseteq M$ (characteristic $0$) with $\rho^{2}=r$, $\sigma^{2}=s$
($r,s \in K$): if an element of the $K$-span of $\{1,\rho,\sigma,\rho\sigma\}$
squares to $x \in K$, then one of $x$, $xr$, $xs$, $xrs$ is a square in $K$.
Degenerate cases ($r$, $s$ a square; $\rho$ or $\sigma$ already in $K$) are not
excluded: the case analysis absorbs them.
\end{lemma}
\begin{proof}\leanok
One quadratic step (\texttt{step}), applied through a master case split on the
square classes of $r$ and $s$.
\end{proof}

\begin{theorem}[Non-squareness in the multiquadratic extension]\label{multiquad_verdicts}\lean{BMThm7SquareClass.Q4_not_square_multiquadratic}\leanok
\uses{sq_transfer, covers_nonsquare}
In every field containing $\Q(a,b)$ together with square roots of $R_{q}$ and
$S_{q}$, neither $Q_{4}$ nor $-2Q_{4}$ is the square of any element of the span
of $\{1,\rho,\sigma,\rho\sigma\}$
(\texttt{Q4\_not\_square\_multiquadratic},
\texttt{n2Q4\_not\_square\_multiquadratic}).
\end{theorem}
\begin{proof}\leanok
The transfer lemma reduces to the eight base-field exclusions of
Theorem~\ref{covers_nonsquare}.
\end{proof}

Since every element of $K(S)$ lies in the span [M: quadratic-tower spanning],
the theorem excludes squares in all of $K(S)$; the kernel boundary is exactly the
span statement, the spanning fact is ordinary mathematics.

\begin{theorem}[$C_{\pm}$-exclusion in the layer]\label{gamma_verdicts}\lean{BMThm7NormCriterion.gamma_not_square_multiquadratic}\leanok
\uses{covers_nonsquare}
For either sign, $\gamma_{\pm}=s_{1n}(s_{1n}\mp\zeta)$ --- an element of the
layer $K(\zeta)$ --- is not the square of any element of the span of
$\{1,\zeta,\sigma,\zeta\sigma\}$ over any field containing $\Q(a,b)$ together
with square roots of $R_{q}$ and $S_{q}$
(\texttt{gamma\_not\_square\_multiquadratic};
supporting kernel lemmas \texttt{step\_over\_layer}, \texttt{layer\_square\_norm},
\texttt{I2\_R2}, \texttt{Rq\_not\_square}).
\end{theorem}
\begin{proof}\leanok
Layer step, coordinate norm, (I2), then \texttt{n2Q4\_not\_square}.
\end{proof}

\begin{proposition}[Geometric integrality of the covers]\label{prop:geomint}
Neither $Q_{4}$ nor $\gamma_{+}$ nor $\gamma_{-}$ is a square in
$\overline{\Q}(S)$; consequently $C_{0}$, $C_{+}$ and $C_{-}$ are dominant,
geometrically integral degree-two covers over the stated open subset of $S$,
away from their branch loci.  ($-2Q_{4}$ is not itself a radicand: it enters as
the layer-norm obstruction class of $\gamma_{\pm}$.)
\end{proposition}
\begin{proof}[Proof {[M]} with {[C]} inputs]
$Q_{4}=F_{1}F_{2}F_{3}$ with the $F_{i}$ distinct irreducible factors, each to
the first power [C: \texttt{terrain\_crosscheck}, original Magma terrain
scripts]; squarefreeness persists over $\overline{\Q}$, so $Q_{4}$ has odd
valuation along every geometric component of each $F_{i}=0$ in
$\overline{\Q}(B)$.  Since $\gcd(F_{i},R_{q}S_{q})=1$ [C], these components are
disjoint from the branch divisors of $S\to B$; ramification indices over them
are $1$, so the valuations stay odd in $\overline{\Q}(S)$ and
$Q_{4}\notin\overline{\Q}(S)^{\times2}$.  The divisor argument applies verbatim
to $-2Q_{4}$ (same zero divisor).  For $C_{\pm}$: the specialised kernel theorem quantifies its coefficients over
$\Q(a,b)$ and is \emph{not} base-changed here.  Instead, the abstract kernel
lemmas \texttt{step\_over\_layer} and \texttt{layer\_square\_norm} --- stated for
an arbitrary characteristic-zero field --- are applied afresh with
$K=\overline{\Q}(a,b)$ (the abstract statements are [K]; this instantiation is
[M]).  Their inputs over $\overline{\Q}(a,b)$ --- $R_{q}$ and $-2Q_{4}$ nonsquare
--- follow from the odd-multiplicity divisor data [C] via
Lemma~\ref{lem:sqfree_basechange}; the assembly reduces geometric squareness of
$\gamma_{\pm}$ to that of $-2Q_{4}$, excluded above [M].  Dominance and generic degree two are
immediate: the radicands are nonzero elements of $K(S)$.
\end{proof}

\begin{lemma}[Negativity certificates at point 1: $Q_4$]\label{Q4_neg}\lean{BMThm7Boundary.Q4_neg}\leanok
\uses{boundary-defs}
$Q_{4}(-51/13,-17/7)<0$.
\end{lemma}
\begin{proof}\leanok
By \texttt{norm\_num}.
\end{proof}

\begin{lemma}[$Q_4 R_q$ at point 1]\label{Q4Rq_neg}\lean{BMThm7Boundary.Q4Rq_neg}\leanok
\uses{boundary-defs}
$Q_{4}(-51/13,-17/7)\cdot R_{q}(-51/13,-17/7)<0$.
\end{lemma}
\begin{proof}\leanok
By \texttt{norm\_num}.
\end{proof}

\begin{lemma}[$Q_4 S_q$ at point 1]\label{Q4Sq_neg}\lean{BMThm7Boundary.Q4Sq_neg}\leanok
\uses{boundary-defs}
$Q_{4}(-51/13,-17/7)\cdot S_{q}(-51/13,-17/7)<0$.
\end{lemma}
\begin{proof}\leanok
By \texttt{norm\_num}.
\end{proof}

\begin{lemma}[$Q_4 R_q S_q$ at point 1]\label{Q4RqSq_neg}\lean{BMThm7Boundary.Q4RqSq_neg}\leanok
\uses{boundary-defs}
$Q_{4}(-51/13,-17/7)\cdot\bigl(R_{q}(-51/13,-17/7)\,S_{q}(-51/13,-17/7)\bigr)<0$.
\end{lemma}
\begin{proof}\leanok
By \texttt{norm\_num}.
\end{proof}

\begin{lemma}[$-2Q_4$ at point 2]\label{n2Q4_neg}\lean{BMThm7Boundary.n2Q4_neg}\leanok
\uses{boundary-defs}
$-2\,Q_{4}(1,3)<0$.
\end{lemma}
\begin{proof}\leanok
By \texttt{norm\_num}.
\end{proof}

\begin{lemma}[$-2Q_4 R_q$ at point 2]\label{n2Q4Rq_neg}\lean{BMThm7Boundary.n2Q4Rq_neg}\leanok
\uses{boundary-defs}
$-2\,Q_{4}(1,3)\cdot R_{q}(1,3)<0$.
\end{lemma}
\begin{proof}\leanok
By \texttt{norm\_num}.
\end{proof}

\begin{lemma}[$-2Q_4 S_q$ at point 2]\label{n2Q4Sq_neg}\lean{BMThm7Boundary.n2Q4Sq_neg}\leanok
\uses{boundary-defs}
$-2\,Q_{4}(1,3)\cdot S_{q}(1,3)<0$.
\end{lemma}
\begin{proof}\leanok
By \texttt{norm\_num}.
\end{proof}

\begin{lemma}[$-2Q_4 R_q S_q$ at point 2]\label{n2Q4RqSq_neg}\lean{BMThm7Boundary.n2Q4RqSq_neg}\leanok
\uses{boundary-defs}
$-2\,Q_{4}(1,3)\cdot\bigl(R_{q}(1,3)\,S_{q}(1,3)\bigr)<0$.
\end{lemma}
\begin{proof}\leanok
By \texttt{norm\_num}.
\end{proof}

\section{What the terrain shows}

The locus inside $S$ where Theorem~7's conclusion is algebraically forced is the proper
curve $S\cap\{Q_{4}=0\}$: six one-dimensional components, each of genus~$0$ [C].  On each
component $R_{q}$ restricts to a square while $S_{q}$ does not; the restriction of $S_{q}$
has constant square class $12=4\cdot 3$ [C], so $w^{2}=S_{q}$ has rational solutions only
where $S_{q}=0$, and every rational zero of $S_{q}$ on each component is a degenerate point of the
chart ($\mathrm{den}=0$ or the parametrisation pole; on $F_{1}$ there are no
rational zeros at all) [C: \texttt{terrain\_crosscheck}], hence inadmissible.  Hence over $\Q$ the conclusion is never automatic: every hypothetical
nondegenerate point of $S(\Q)$ must lift through a nontrivial double cover, and the
Conjecture-1-refutation scenario via these curves is closed (no admissible points) [C+M].

Away from the degeneracy locus each of the three disjuncts demands lifting through a
double cover that is nontrivial over the splitting extension $K(S)$, so the
liftable rational points of each cover form a thin subset of type II [M] ---
thinness here is essentially definitional (the image of the rational points of a
dominant degree-two cover); what Hilbert irreducibility controls is the
\emph{complement}: outside a thin set, specialisation preserves irreducibility.
(Irreducibility over $K(S)$: Theorem~\ref{multiquad_verdicts} for $C_{0}$ and
Theorem~\ref{gamma_verdicts} for $C_{\pm}$; dominance, degree two, and geometric
integrality away from the branch locus are Proposition~\ref{prop:geomint}.)
Thinness bounds what a generic-point argument can see, not what
rational points exist; we found no geometric mechanism that decides the lifting, and no
descent route achieving finiteness (the naive twist family is unbounded, and none of the
standard finiteness devices --- ramification restriction, local conditions, Selmer-type
bounds --- was found by us to apply); and the field-dependence begins
already at Buchholz--MacDougall's own excluded case, over the quadratic extension suggested by
the square class, $\Q(\sqrt{3})$ ($12=4\cdot 3$), while Stroeker's points live over $\Q(\sqrt{3441})$.

\begin{proposition}[Normalization bridge]\label{prop:normalization_bridge}\lean{BMThm7Terrain.conjecture1_iff_terrain}\leanok
There exists a $\Q$-derived quartic with four distinct rational roots (the
$\{0,1,a,b\}$ normal form of Conjecture~1) if and only if there exists a
nondegenerate rational point of the splitting terrain: $(a,b,z,w)\in\Q^{4}$ with
$z^{2}=R_{q}(a,b)$, $w^{2}=S_{q}(a,b)$, $\mathrm{den}(a,b)\neq0$, and the four
roots $\{1,a,b,c\}$, $c=-ab/\mathrm{den}$, pairwise distinct and nonzero.
Equivalently: Conjecture~1 holds iff $S(\Q)$ has no nondegenerate point.
\end{proposition}
\begin{proof}[Proof; both directions, every normalisation condition tracked]\leanok
Kernel-verified as of 2026-08-02: \texttt{Gap/BMThm7Terrain.lean} proves
\texttt{conjecture1\_iff\_terrain} (with \texttt{kderived\_of\_terrain},
\texttt{terrain\_of\_kderived}, the three K-derivedness transport lemmas, and the
quadratic splitting criteria) --- all standard-axiom.  The readable account:
($\Leftarrow$, terrain point $\to$ quartic)  Set
$f=(x-1)(x-a)(x-b)(x-c)$, monic; $e_{3}=c\,\mathrm{den}+ab=0$ by construction, so
$f'(0)=-e_{3}=0$ and $f'=x\,(4x^{2}-3e_{1}x+2e_{2})$.  With
$e_{1}=s_{1n}/\mathrm{den}$, $e_{2}=s_{2n}/\mathrm{den}$ one computes
$9e_{1}^{2}-32e_{2}=R_{q}/\mathrm{den}^{2}=(z/\mathrm{den})^{2}$, so the quadratic
factor splits with explicit roots $(3e_{1}\pm z/\mathrm{den})/8$; likewise
$f''=12x^{2}-6e_{1}x+2e_{2}$ has $9e_{1}^{2}-24e_{2}=S_{q}/\mathrm{den}^{2}
=(w/\mathrm{den})^{2}$ and splits; $f'''$ is linear.  So $f$ is $\Q$-derived with
four distinct rational roots; translating the root $1$ to $0$ and scaling a
second (nonzero, since roots are distinct and nonzero) root to $1$ is an affine
change preserving $\Q$-derivedness, yielding the $\{0,1,a',b'\}$ normal form with
distinct parameters.
($\Rightarrow$, quartic $\to$ terrain point)  In normal form
$f=x(x-1)(x-a)(x-b)$, $\Q$-derived, roots distinct.  $f'$ splits, so it has a
rational root $x_{0}$; distinct roots make $f$ squarefree, and at a root $r$ of
$f$, $f'(r)=\prod_{j\ne i}(r-r_{j})\neq0$, so $x_{0}$ is \emph{not} a root of
$f$.  Translate $x\mapsto x+x_{0}$: the translated roots
$\{-x_{0},1-x_{0},a-x_{0},b-x_{0}\}$ are all nonzero, and $f'(0)=0$ gives
$e_{3}=0$.  Scale by the nonzero root $-x_{0}$, sending it to $1$; scaling
multiplies $e_{3}$ by a nonzero cube, so $e_{3}=0$ persists.  Writing the
resulting roots $\{1,A,B,C\}$: $e_{3}=C\,\mathrm{den}(A,B)+AB=0$; if
$\mathrm{den}(A,B)=0$ then $AB=0$, contradicting nonzero roots --- so
$\mathrm{den}(A,B)\neq0$ automatically, and $C=-AB/\mathrm{den}$.  The split
quadratic factors of $f'$ and $f''$ have rational root differences with
$(r_{+}-r_{-})^{2}=R_{q}/(16\,\mathrm{den}^{2})$ and
$(u_{+}-u_{-})^{2}=S_{q}/(36\,\mathrm{den}^{2})$; set
$z:=4\,\mathrm{den}\,(r_{+}-r_{-})$, $w:=6\,\mathrm{den}\,(u_{+}-u_{-})$, so
$z^{2}=R_{q}(A,B)$ and $w^{2}=S_{q}(A,B)$.  All conditions ($\mathrm{den}\neq0$;
distinctness; nonzeroness; the $z,w$ witnesses) are tracked; the choices of
critical point and scaling root are immaterial, as only existence is asserted.
\end{proof}

The reading is an assessment [A] anchored by machine artefacts, not a theorem: within the
proof architectures we surveyed, each escape hatch for a Theorem-7 proof avoiding
Conjecture~1 is shut at the level computation can shut it, and what remains open is that
$S(\Q)$ has no nondegenerate points --- equivalent to Conjecture~1 by
Proposition~\ref{prop:normalization_bridge}.  Formally, sufficiency
(Theorem~\ref{thm7_of_conjecture1}) and the normalization bridge
(Proposition~\ref{prop:normalization_bridge}) are both kernel theorems; necessity of
Conjecture~1 for Theorem~7, and the nonexistence of routes unknown to us, are not
claimed.  Every route to Theorem~7 known to us passes through the
white node.
